\documentclass[10pt,a4paper]{amsart}
\usepackage[margin=19mm]{geometry}
\usepackage{amsmath,amssymb,mathtools}
\usepackage[scr=rsfs,cal=euler]{mathalfa}
\usepackage{xfrac}
\usepackage[unicode,hidelinks,bookmarks=false]{hyperref}
\newcommand{\ie}{i.e.\ }
\newcommand{\cf}{cf.\ }
\newcommand{\resp}{respectively\ }
\newcommand{\Subsection}{\subsection}
\providecommand{\nobreakdash}{\nobreak\hbox{-}}
\setlength{\emergencystretch}{2em}
\multlinegap0pt
\usepackage{bm}
\usepackage{bbm}
\usepackage{dsfont}
\usepackage{xspace}
\usepackage{cancel}
\usepackage{mathtools}
\usepackage{amsthm}
\makeatletter
\@ifundefined{theorem}{\newtheorem{theorem}{Theorem}}{}
\@ifundefined{lemma}{\newtheorem{lemma}[theorem]{Lemma}}{}
\makeatother
\newcommand{\lcm}{\textup{lcm}}
\renewcommand{\mod}[2]{\equiv#1\textup{ (mod }#2\textup{)}}
\title[Covering systems with odd moduli]{Covering systems with odd moduli}
\author{Joshua Harrington, Yewen Sun, Wing Hong Tony Wong}
\date{Source: arXiv:2104.00602v1 (CC BY 4.0)}
\begin{document}
\maketitle

\noindent\emph{Source and licence.} Covering systems with odd moduli. Version arXiv:2104.00602v1 (CC BY 4.0), distributed under the Creative Commons Attribution 4.0 International licence, as recorded in the arXiv licence metadata for this exact version and shown on the abstract page https://arxiv.org/abs/2104.00602v1. Full rights notice: licenses/arxiv-2104.00602-CC-BY-4.0.txt.\par\smallskip

\noindent\textit{Editorial scope.} Counted unit: the Theorem whose source label is \texttt{thm:3sqrfreeimpliesodd} in the article (source0.tex lines 354--356, source label \texttt{thm:3sqrfreeimpliesodd}), delivered together with the article's own complete original proof of it (source0.tex lines 358--403, one \texttt{proof} environment of 46 lines). No mathematics is written, repaired or re-derived: statement and proof are the article's own text line for line. Required context retained uncounted: lem:shifted (lines 342--345). Every definition, display and label of those lines is retained unchanged. Omitted from this package and not redistributed: 1--341, 346--353, 357--357, 404--1135. Nothing inside a retained excerpt is omitted or altered: every retained body is delivered exactly as the article's lines are written, so no omission receipt of any kind accompanies this package. Citations: the retained excerpts carry no citation command at all, so no bibliography entry of the article is transcribed and no reference is redistributed. Reference adaptation: none was needed and none was made; every reference inside the delivered unit resolves inside it, so no unresolved reference or citation is produced. Printed slips kept literal: no printed word, symbol, hypothesis or number of the counted statement or of its proof was altered. Layout and class: the article's own class is \texttt{article} with packages including amsmath, amsthm, amsfonts, amssymb, fullpage, multirow, array, bbm, authblk, verbatim, tikz, float, enumerate, diagbox; the wrapper is the lane's pinned \texttt{amsart} editorial wrapper at 10pt on A4, which restates the theorem-like environments the retained text needs and adds the article's own macro definitions that the retained text needs (\texttt{\textbackslash lcm}, \texttt{\textbackslash mod}), with extra wrapper packages (bm, bbm, dsfont, xspace, cancel, mathtools). The delivered numbering is the unit's own. Assets: no retained span contains a figure environment, a graphics-inclusion command, a PDF-page inclusion, a table or a TikZ picture; no image file is copied into this package, the package declares no asset dependency and no image file is redistributed with it. Licence: version arXiv:2104.00602v1 of this article is distributed under the Creative Commons Attribution 4.0 International licence, as recorded in the arXiv licence metadata for this exact version and shown on the abstract page https://arxiv.org/abs/2104.00602v1. A full rights notice is delivered at licenses/arxiv-2104.00602-CC-BY-4.0.txt. Characters: every retained excerpt is pure ASCII, so the delivered engine, XeLaTeX with its default Latin Modern text font, reports no missing character.
\begin{lemma}\label{lem:shifted}
Let $d$ be an integer.  Given a covering system $\mathcal{C}=\{\mod{r_j}{m_j}: 1\leq j\leq k\}$ of $S\subseteq\mathbb{Z}$, the set $\mathcal{C}_
d=\{r_j+d\pmod{m_j}:1\leq j\leq k\}$ is a covering of the set $S_d=\{s+d:s\in S\}$.
\end{lemma}

\begin{theorem}\label{thm:3sqrfreeimpliesodd}
Let $p\geq3$ be a prime. If there exists a covering system of the integers such that all moduli are odd, square-free, and distinct except that $p$ is used exactly twice as a modulus, then there exists an odd covering of the integers.
\end{theorem}

\begin{proof}
Suppose that $\mathcal{C}_0$ is a covering system of the integers such that all moduli are odd, square-free, and distinct except that $p$ is used exactly twice as a modulus. By Lemma~\ref{lem:shifted}, we may assume that $\mod{0}{p}$ and $\mod{1}{p}$ are congruences in $\mathcal{C}_0$. Let $k$ be a nonnegative integer and let $\ell_\xi$ be nonnegative integers for each $2\leq\xi\leq p-1$ such that
\begin{align*}
\mathcal{C}_0&=\{\mod{0}{p},\mod{1}{p}\}\cup\{\mod{r_j}{m_j}:1\leq j\leq k\}\cup\\
&\hspace{20pt}\bigcup_{\xi=2}^{p-1}\{\mod{r_{\xi,j}}{pm_{\xi,j}}:1\leq j\leq\ell_\xi\},
\end{align*}
where $p\nmid m_j$ for all $1\leq j\leq k$, and $p\nmid m_{\xi,j}$ and $r_{\xi,j}\equiv\mod{\xi}{p}$ for all $1\leq j\leq k_\xi$ and $2\leq\xi\leq p-1$.

We claim that for each $2\leq\xi\leq p-1$, the set of congruences 
$$\mathcal{C}_\xi=\{\mod{r_j}{m_j}:1\leq j\leq k\}\cup\{\mod{r_{\xi,j}}{m_{\xi,j}}:1\leq j\leq\ell_\xi\}$$
forms a covering system of the integers. To see this, consider any integer $r$. Let
$$m=\lcm\big(\{m_j:1\leq j\leq k\}\cup\{m_{\xi,j}:1\leq j\leq\ell_\xi\}\big)$$
and let $n$ be an integer that satisfies the congruence system
$$\begin{cases}
n\equiv\mod{\xi}{p},\\
n\equiv\mod{r}{m}.
\end{cases}$$
Since $\mathcal{C}_0$ is a covering system of the integers, $n$ must satisfy one of the congruences in
$$
\{\mod{r_j}{m_j}:1\leq j\leq k\}\cup\{\mod{r_{\xi,j}}{pm_{\xi,j}}:1\leq j\leq\ell_\xi\}.$$
As a result, if $n\equiv\mod{r_{\tilde{j}}}{m_{\tilde{j}}}$ for some $1\leq\tilde{j}\leq k$, then
$$r\equiv n\equiv\mod{r_{\tilde{j}}}{m_{\tilde{j}}};$$
if $n\equiv\mod{r_{\xi,\tilde{j}}}{pm_{\xi,\tilde{j}}}$ for some $1\leq\tilde{j}\leq\ell_\xi$, then
$$r\equiv n\equiv\mod{r_{\xi,\tilde{j}}}{m_{\xi,\tilde{j}}}.$$
This completes the proof of our claim.

Let $q$ be an odd prime such that $q\nmid pm$. For each $1\leq i\leq q-2$, $2\leq\xi\leq p-2$, and $1\leq j\leq\ell_\xi$, let $r_{\xi,i,j}$ be an integer that satisfies the congruence system
$$\begin{cases}
r_{\xi,i,j}\equiv\mod{\xi\cdot p^i}{p^{i+1}},\\
r_{\xi,i,j}\equiv\mod{r_{\xi,j}}{m_{\xi,j}}.
\end{cases}$$
Moreover, for each $0\leq i\leq q-1$, let $s_i$ be an integer that satisfies the congruence system
$$\begin{cases}
s_i\equiv\mod{0}{p^i},\\
s_i\equiv\mod{i}{q}.
\end{cases}$$
Now, we are ready to build an odd covering system based on the existence of $\mathcal{C}_0$. Consider
\begin{align*}
\mathcal{C}&=\{\mod{p^i}{p^{i+1}}:0\leq i\leq q-2\}\cup\{\mod{s_i}{p^iq}:0\leq i\leq q-1\}\\
&\hspace{20pt}\cup\{\mod{r_j}{m_j}:1\leq j\leq k\}\cup\bigcup_{\xi=2}^{p-1}\{\mod{r_{\xi,i,j}}{p^{i+1}m_{\xi,j}}:0\leq i\leq q-2,1\leq j\leq\ell_\xi\}.
\end{align*}

It is easy to see that all moduli of $\mathcal{C}$ are odd and distinct. To show that $\mathcal{C}$ is a covering system of the integers, let $N$ be an arbitrary integer, and let $i_0$ be the maximum integer such that $p^{i_0}\mid N$. If $0\leq i_0\leq q-2$, then either $N\equiv\mod{p^{i_0}}{p^{i_0+1}}$, which is trivially covered by $\mathcal{C}$, or $N\equiv\mod{\xi\cdot p^{i_0}}{p^{i_0+1}}$ for some $2\leq\xi\leq p-1$, which is covered by 
$$\{\mod{r_j}{m_j}:1\leq j\leq k\}\cup\{\mod{r_{\xi,i_0,j}}{p^{i_0+1}m_{\xi,j}}:1\leq j\leq\ell_\xi\}$$
since $C_\xi$ is a covering system of the integers. Lastly, if $i_0\geq q-1$, then $N$ is covered by $\{\mod{s_i}{p^iq}:0\leq i\leq q-1\}$.
\end{proof}

\end{document}
